% Draft replacement for the implementation/protocol/experiments part of the paper.
% The text is intended to be pasted after the mathematical and algorithmic sections.
\providecommand{\ratioiqr}[3]{\shortstack[c]{#1\\[-0.25ex]{\tiny #2--#3}}}

\section{Implementation}
\label{sec:implementation}

We implemented the Morse persistence pipeline as a C++ prototype with a pure C++ API and
a Python interface.  The C++ layer contains the complex representation, the Morse-sequence
constructors, the reference-map construction, and the persistence reducer.  The Python
interface is used to generate benchmark families, inspect Morse sequences and critical
complexes, run reproducible experiments, and compare with external persistence
implementations.  A GUDHI-facing C++ interface, starting directly from a
\texttt{Gudhi::Simplex\_tree}, is under development, but it is not used as an experimental
baseline in the timings reported below.

The pipeline has the following stages:
\[
  (K,F)
  \longrightarrow
  \hbox{Morse sequence}
  \longrightarrow
  \hbox{reference annotations}
  \longrightarrow
  \hbox{persistence on the critical complex}.
\]
The complex is a finite filtered simplicial complex.  After construction, each simplex has
a compact integer identifier, a dimension, a scalar filtration value, a compressed
filtration level, its codimension-one boundary, its coboundary, and its position in a
filtration-compatible order.  The current prototype materializes the complex explicitly;
it is therefore not meant to compete with specialized implicit Vietoris--Rips
constructors.  This choice is deliberate: it gives deterministic boundary and coboundary
access, which is useful for validating the reference recurrence and for comparing
Morse-sequence strategies under the same reducer.

The Morse-sequence constructor is a strategy parameter.  The current implementation
contains the saturated constructor, the same-level reduction constructor, the
paper-faithful \texttt{f-max} and \texttt{f-min} constructors, a plateau-greedy variant,
and several flooding variants.  All strategies return the same object: a sequence of
critical insertions and same-level regular pairs, the list of critical simplices, and the
pairing map.  The persistence reducer is independent of the chosen strategy.

For a fixed Morse sequence, the reference map expresses each simplex as a linear
combination of critical simplices.  Over $\mathbb{F}_2$, annotations are sparse sorted
sets of critical labels and addition is symmetric difference.  Over an odd prime field
$\mathbb{F}_p$, annotations are sparse sorted vectors of
$(\hbox{critical label},\hbox{coefficient})$ pairs, and oriented boundary signs are
reduced modulo $p$.  Composite rings $\mathbb{Z}_n$ are not used in the barcode API,
because the current reduction assumes field coefficients.

The reducer stores inverse annotation lists: for each critical label, it keeps the list of
annotations currently containing that label.  Pivot elimination then touches only
annotations that contain the pivot label.  The implementation records the initial
annotation volume, maximum annotation size, inverse-list entries, update counts, and the
time spent in sequence construction, reference construction, and reduction.  The ongoing
GUDHI integration follows the same internal API: it will build the incidence and order
data required by these routines from a \texttt{Gudhi::Simplex\_tree}, while keeping the
reducer and Morse strategies unchanged.

\section{Experimental protocol}
\label{sec:experimental-protocol}

All reported experiments check barcode correctness against ordinary persistence on the
full filtered complex.  Zero-length intervals inside one filtration level are ignored in
the comparison, since different valid reductions may pair simplices differently inside a
plateau without changing the persistence module.  The test suite also checks sequence
validity, the reference recurrence, the coreference recurrence, prime-field reductions,
and agreement between the compact C++ and Python interfaces.

\paragraph{Benchmark families.}
The experiments use three types of data.  First, synthetic lower-star, plateau, and random
Rips complexes are used to test behavior under controlled changes of size, density, and
plateau structure.  Second, triangulated terrain and image-grid examples provide
structured two-dimensional filtrations with many equal values.  Third, selected
Vietoris--Rips complexes from the persistent-homology roadmap benchmark collection are
used as external reference instances.

\paragraph{Why GUDHI is the external baseline.}
We use GUDHI as the main external baseline because the Roadmap paper
\cite{OtterPorterTillmannGrindrodHarrington2017} compared several persistent-homology
software packages on a common set of data sets and includes GUDHI among the reference
implementations.  This is useful for calibration: by comparing our code with GUDHI on
Roadmap-style inputs, the reader can relate our timings to a software package that has
already been positioned in a broader benchmark study.  We do not claim that GUDHI is the
fastest persistent-homology implementation in all settings.  For Vietoris--Rips
filtrations, specialized Ripser-derived implementations are often faster
\cite{BauerRipser2021,ZhangXiaoWangRipserPlusPlus2020,BurellaPerezHaukeLupoCaorsiDassattiGiottoPH2021}.
GUDHI is chosen here because it is mature, widely used, Roadmap-calibrated, and a natural
C++ integration target.

The comparison reported below starts from our explicit filtered complex and times the
compact C++ Morse backend against GUDHI persistence on the same list of simplices and
filtration values.  The GUDHI \texttt{SimplexTree} is populated before the timed
persistence call; thus the GUDHI timing is favorable to GUDHI and does not include
construction from the simplex list.  The separate GUDHI-facing C++ interface is still in
development and is therefore not used for the present paper-facing timing comparison.

For all timing tables, a ratio larger than one means that the denominator method is faster
than the numerator baseline.  Thus GUDHI/Morse larger than one means that the Morse
pipeline is faster than GUDHI.
When a table aggregates several paired observations, ratio columns report the median with
\(Q_1\)--\(Q_3\) underneath.  In the Roadmap portfolio table these intervals describe
case-to-case variation over the selected objects.  Rows with one aggregate observation are
marked as point estimates.

\section{Experimental results}
\label{sec:experimental-results}

\subsection{Compact C++ backend versus GUDHI}

Table~\ref{tab:paper-roadmap-selected} reports the backend comparison on three selected
Roadmap Rips complexes.  The same filtered complex is used for our C++ Morse backend and
for GUDHI.  We show four representative Morse-sequence strategies:
\texttt{saturated}, \texttt{f-max}, \texttt{f-min}, and \texttt{same-level}.

\begin{table*}[t]
\centering
\scriptsize
\setlength{\tabcolsep}{2.5pt}
\begin{tabular}{llrrrrr}
\toprule
Dataset & Strategy & Simplices & Critical \% & Morse ($\mu$s) & GUDHI ($\mu$s) & GUDHI/Morse (point) \\
\midrule
\texttt{random50-16d} & \texttt{saturated} & 20,875 & 89.1 & 6,731 & 9,298 & 1.38 \\
\texttt{random50-16d} & \texttt{f-max} & 20,875 & 89.1 & 5,735 & 9,298 & 1.62 \\
\texttt{random50-16d} & \texttt{f-min} & 20,875 & 89.1 & 6,039 & 9,298 & 1.54 \\
\texttt{random50-16d} & \texttt{same-level} & 20,875 & 99.3 & 5,431 & 9,298 & 1.71 \\
\midrule
\texttt{random100-4d} & \texttt{saturated} & 166,750 & 94.2 & 91,902 & 82,858 & 0.90 \\
\texttt{random100-4d} & \texttt{f-max} & 166,750 & 94.2 & 82,689 & 82,858 & 1.00 \\
\texttt{random100-4d} & \texttt{f-min} & 166,750 & 94.2 & 89,447 & 82,858 & 0.93 \\
\texttt{random100-4d} & \texttt{same-level} & 166,750 & 99.8 & 76,253 & 82,858 & 1.09 \\
\midrule
\texttt{senate104} & \texttt{saturated} & 182,207 & 94.4 & 91,024 & 87,125 & 0.96 \\
\texttt{senate104} & \texttt{f-max} & 182,207 & 94.4 & 80,962 & 87,125 & 1.08 \\
\texttt{senate104} & \texttt{f-min} & 182,207 & 94.4 & 92,470 & 87,125 & 0.94 \\
\texttt{senate104} & \texttt{same-level} & 182,207 & 99.9 & 80,891 & 87,125 & 1.08 \\
\bottomrule
\end{tabular}
\caption{Backend comparison on selected Roadmap Rips complexes.  GUDHI is timed on the
same filtered complexes, after construction of its \texttt{SimplexTree}.  Each row is one
aggregate observation, so the ratio is a point estimate.}
\label{tab:paper-roadmap-selected}
\end{table*}

The selected Roadmap examples support two observations.  First, the compact C++ Morse
backend is faster than GUDHI on most of these rows, with the strongest portfolio rows
coming from \texttt{same-level} and \texttt{f-max}.  Some case-level intervals still
straddle one, so this should be read as targeted evidence on selected Roadmap objects, not
as a universal speedup claim.  Second, the number of critical cells is not a sufficient predictor of time:
\texttt{same-level} leaves nearly every simplex critical, but it is the fastest average
strategy on these Roadmap instances because it performs much less annotation mutation in
the reducer.

\begin{table*}[t]
\centering
\scriptsize
\setlength{\tabcolsep}{3pt}
\begin{tabular}{lrrcc}
\toprule
Strategy & Critical \% & Morse ($\mu$s) & Standard/Morse & GUDHI/Morse \\
\midrule
\texttt{same-level}      & 99.8 & 54,192 & \ratioiqr{4.74}{3.23}{6.40} & \ratioiqr{1.09}{1.08}{1.40} \\
\texttt{f-max}           & 94.0 & 56,462 & \ratioiqr{4.57}{3.16}{6.01} & \ratioiqr{1.08}{1.04}{1.35} \\
\texttt{f-min}           & 94.0 & 62,652 & \ratioiqr{4.17}{2.84}{5.51} & \ratioiqr{0.94}{0.93}{1.24} \\
\texttt{saturated}       & 94.0 & 63,219 & \ratioiqr{3.80}{2.66}{5.22} & \ratioiqr{0.96}{0.93}{1.17} \\
\texttt{plateau-greedy}  & 94.0 & 74,674 & \ratioiqr{2.95}{2.12}{4.33} & \ratioiqr{0.81}{0.79}{0.92} \\
\texttt{flooding-min}    & 94.0 & 75,659 & \ratioiqr{2.58}{1.90}{4.25} & \ratioiqr{0.80}{0.78}{0.85} \\
\texttt{flooding-max}    & 94.0 & 77,831 & \ratioiqr{2.69}{1.93}{4.16} & \ratioiqr{0.78}{0.76}{0.85} \\
\texttt{flooding-maxmin} & 94.0 & 79,268 & \ratioiqr{2.50}{1.85}{3.97} & \ratioiqr{0.75}{0.75}{0.84} \\
\texttt{flooding-minmax} & 94.0 & 84,056 & \ratioiqr{2.66}{1.86}{3.91} & \ratioiqr{0.71}{0.70}{0.82} \\
\bottomrule
\end{tabular}
\caption{Average Roadmap timings over the three selected Roadmap objects.  The first four
strategies are the main same-level Morse-sequence constructors; the flooding variants are
included as a diagnostic portfolio.  Ratio entries are case-level medians with
\(Q_1\)--\(Q_3\) underneath over the three objects.}
\label{tab:paper-roadmap-portfolio}
\end{table*}

\subsection{Structural measurements}

The complexity estimates in Section~\ref{sec:complexity} predict that running time and
memory should be controlled not only by the number of critical simplices, but also by the
reference working set and by annotation volume.  Table~\ref{tab:paper-structural-roadmap}
therefore reports the corresponding structural quantities for the four main strategies on
the selected Roadmap objects.  Here \(n_\partial^+=|\Wbd|\), and \(S_\partial\) is measured
as the initial reducer annotation volume after fused reference construction.  In these
\(\mathbb{F}_2\) runs, the number of inverse-list entries at reducer setup is equal to
\(S_\partial\), since each stored critical label contributes one inverse-list entry.

\begin{table*}[t]
\centering
\small
\begin{tabular}{lrrrrrr}
\toprule
Strategy & Avg. \(|K|\) & \(c/|K|\) & \(n_\partial^+/|K|\)
  & \(S_\partial/|K|\) & Face scans/\(|K|\) & XOR input/\(|K|\) \\
\midrule
\texttt{same-level} & 123,277 & 99.7\% & 99.8\% & 1.00 & 2.95 & 0.086 \\
\texttt{f-min}      & 123,277 & 92.6\% & 96.3\% & 1.22 & 2.77 & 0.354 \\
\texttt{f-max}      & 123,277 & 92.6\% & 96.3\% & 1.27 & 2.77 & 0.063 \\
\texttt{saturated}  & 123,277 & 92.6\% & 96.3\% & 1.22 & 2.77 & 0.354 \\
\bottomrule
\end{tabular}
\caption{Reference-side structural measurements on the selected Roadmap Rips objects.
Values are row means over the three selected objects.  Face scans count planned boundary
queries for critical simplices.  XOR input is the total size of annotation inputs touched
by pivot-elimination updates, normalized by \(|K|\).}
\label{tab:paper-structural-roadmap}
\end{table*}

The table explains why critical-cell count alone is not a reliable predictor.  The
\texttt{same-level} sequence leaves almost every simplex critical, but its annotations are
close to unit size and its elimination updates touch little data.  Conversely,
\texttt{f-min} and \texttt{saturated} have fewer critical simplices, but larger update
volume on these examples.  The dual quantities \(n_\delta^+=|\Wcobd|\) and \(S_\delta\)
are instrumented in the implementation and are used in the coreference diagnostics, but
they are not part of the main timing comparison because the paper-facing backend uses the
reference direction.  In the dense Rips examples tested so far, the naive coboundary-based
coreference recurrence has substantially larger update volume than the reference
recurrence.

\subsection{Coefficient fields}

The timing tables above use $\mathbb{F}_2$, which is the coefficient field used in the
GUDHI comparison.  We also ran a prime-field validation over $\mathbb{F}_3$ and
$\mathbb{F}_5$ on synthetic lower-star, plateau, and Rips complexes, with the
\texttt{saturated}, \texttt{f-max}, \texttt{f-min}, and \texttt{same-level} strategies.
All prime-field Morse barcodes matched the corresponding standard prime-field reducer.

\begin{table*}[t]
\centering
\small
\setlength{\tabcolsep}{3pt}
\begin{tabular}{lrccc}
\toprule
Field & Cases & $\mathbb{F}_p/\mathbb{F}_2$ Morse
      & $\mathbb{F}_p/\mathbb{F}_2$ standard
      & Morse speedup vs standard \\
\midrule
$\mathbb{F}_3$ & 180 & \ratioiqr{1.231}{1.190}{1.267} & \ratioiqr{1.040}{1.022}{1.066} & \ratioiqr{1.466}{1.078}{1.898} \\
$\mathbb{F}_5$ & 180 & \ratioiqr{1.216}{1.175}{1.253} & \ratioiqr{1.041}{1.023}{1.062} & \ratioiqr{1.465}{1.122}{1.969} \\
\bottomrule
\end{tabular}
\caption{Prime-field validation and overhead check.  Validation mismatches: 0.  Morse
timings exclude sequence construction, which is coefficient-independent and reported
separately in the CSV artifacts.  Ratio entries are medians with \(Q_1\)--\(Q_3\)
underneath over the aggregate benchmark rows.}
\label{tab:paper-prime-field}
\end{table*}

\paragraph{Interpretation.}
The experiments support the following conservative conclusions.  The Morse-frame pipeline
reproduces ordinary persistence on all tested instances.  On the selected Roadmap Rips
objects, the compact C++ backend is faster than GUDHI on average, even though the GUDHI
timing excludes simplex-tree construction.  This is the performance comparison that is
most relevant to the algorithmic backend at this stage.  The GUDHI-facing C++ interface
remains an integration target, but it should be evaluated separately once the interface
path has been stabilized and optimized.  Finally, the Roadmap/GUDHI comparison should be
read as a calibrated baseline, not as a claim that the prototype has been compared
directly with every package in the Roadmap paper or with specialized Rips solvers.
